%
% Copyright (C) 1995, 2000, 2013 American Mathematical Society.
% Copyright (C) 2016-2025 LaTeX Project and American Mathematical Society.
%
% This work may be distributed and/or modified under the
% conditions of the LaTeX Project Public License, either version 1.3c
% of this license or (at your option) any later version.
% The latest version of this license is in
%   https://www.latex-project.org/lppl.txt
% and version 1.3c or later is part of all distributions of LaTeX
% version 2005/12/01 or later.
%
% This work has the LPPL maintenance status `maintained'.
%
% The Current Maintainer of this work is the LaTeX Project.
%
\documentclass[leqno,titlepage,openany]{amsldoc}[1999/12/13]

\def\MaintainedByLaTeXTeam#1{%
\begin{center}%
\fbox{\fbox{\begin{tabular}{@{}l@{}}%
This file is maintained by the \LaTeX{} Project team.\\%
Bug reports can be opened (category \texttt{#1}) at\\%
\url{https://latex-project.org/bugs/}.\end{tabular}}}\end{center}}

\ifx\UndEfiNed\url
  \ClassError{amsldoc}{%
    This version of amsldoc.tex must be processed\MessageBreak
    with a newer version of amsldoc.cls (2.02 or later).}{}
\fi

\renewcommand\bibname{References}

%    Definitions that should be moved to amsldoc.cls
%    when it is next updated
\makeatletter
\let\cleardouble@page\cleardoublepage
\AtBeginDocument{%
  \ifx\cleardouble@page\cleardoublepage
    \def\cleardoublepage{\clearpage{\pagestyle{empty}\cleardouble@page}}
  \fi
}

% #1 is empty or one of class, environment, option, package
% could use different font in each case
\makeatletter
\def\@category@index#1#2{%
    \autoindex{%
        #2\actualchar\string\index#1font{#2}%
        \ifx\@nil#1\@nil\else\space#1\fi
    }%
}

\let\indexpackagefont\textsf
\let\indexclassfont\textsf
\let\indexoptionfont\texttt
\let\indexenvironmentfont\texttt
\let\indexfont\texttt
\makeatother
%    For a package that won't be indexed, i.e. amsmath
\newcommand{\nipkg}{\textsf}
\def\tikz/{Ti\textit{k}Z}
%    End definitions for amsldoc.cls

\title{User's Guide for the \nipkg{amsmath} Package (Version~2.1)}
\author{American Mathematical Society, \LaTeX\ Project}
\date{1999-12-13\\(revised 2002-02-25, 2016-11-14, 2018-04-05, 2019-10-14, 2020-02-18)}
\makeatletter
\def\@thanks{\bigskip\MaintainedByLaTeXTeam{amslatex}}
\makeatother

% Very minor changes are not listed in the \date above:
%
% 2021-08-22: reported typo fixed mathools -> mathtools
% 2025-05-23: drop spurious option list


%    Use the amsmath package and amscd package in order to print
%    examples.
\usepackage{amsmath}
\usepackage{amscd}

\usepackage{imakeidx}
%    Need alternate coding for vertical bar to get it into index.
%    Based on \qcbang and \cnbang in amsldoc.cls, refined by David Carlisle.
%    Give it a sort field that will be at the end of the symbols;
%    must avoid " and | which have special meanings.
\DeclareRobustCommand{\qcvert}{\qc\|\index{{z1}@{\ntt\char'174}}}
\DeclareRobustCommand{\cnvert}{%
  \ncn{\|}\index{{z0}@{\ntt\bslchar\char'174}}}
%    Need alternate index sort field for bang to allow & to sort first
\DeclareRobustCommand{\altcnbang}{%
  \ncn{\!}\index{"!b@{\ntt\bslchar\qcbang}}}
\DeclareRobustCommand{\cnamp}{%
  \index{"!a@\texttt{\&}}}

\makeindex[intoc] % generate index data
\providecommand{\see}[2]{\textit{see} #1}

%    The amsldoc class includes a number of features useful for
%    documentation about TeX, including:
%
%    ---Commands \tex/, \amstex/, \latex/, ... for uniform treatment
%    of the various logos and easy handling of following spaces.
%
%    ---Commands for printing various common elements: \cn for command
%    names, \fn for file names (including font-file names), \env for
%    environments, \pkg and \cls for packages and classes, etc.

%    Many of the command names used here are rather long and will
%    contribute to poor linebreaking if we follow the \latex/ practice
%    of not hyphenating anything set in tt font; instead we selectively
%    allow some hyphenation.
\allowtthyphens % defined in amsldoc.cls

\hyphenation{ac-cent-ed-sym-bol add-to-counter add-to-length align-at
  aligned-at allow-dis-play-breaks ams-art ams-cd ams-la-tex amsl-doc
  ams-symb ams-tex ams-text ams-xtra bmatrix bold-sym-bol cen-ter-tags
  eqn-ar-ray idots-int int-lim-its latex med-space neg-med-space
  neg-thick-space neg-thin-space no-int-lim-its no-name-lim-its
  over-left-arrow over-left-right-arrow over-right-arrow pmatrix
  qed-sym-bol set-length side-set small-er tbinom the-equa-tion
  thick-space thin-space un-der-left-arrow un-der-left-right-arrow
  un-der-right-arrow use-pack-age var-inj-lim var-proj-lim vmatrix
  xalign-at xx-align-at}

%    Prepare for illustrating the \vec example
\newcommand{\vect}[1]{\mathbf{#1}}

\newcommand{\booktitle}[1]{\textit{#1}}
\newcommand{\journalname}[1]{\textit{#1}}
\newcommand{\seriesname}[1]{\textit{#1}}

%    Command to insert and index a particular phrase. Doesn't work for
%    certain kinds of special characters in the argument.
\newcommand{\ii}[1]{#1\index{#1}}

\newcommand{\vstrut}[1]{\vrule width0pt height#1\relax}

%    An environment for presenting comprehensive address information:
\newenvironment{infoaddress}{%
  \par\topsep\medskipamount
  \trivlist\centering
  \item[]%
  \begin{minipage}{.7\columnwidth}%
  \raggedright
}{%
  \end{minipage}%
  \endtrivlist
}

\newenvironment{eqxample}{%
  \par\addvspace\medskipamount
  \noindent\begin{minipage}{.5\columnwidth}%
  \def\producing{\end{minipage}\begin{minipage}{.5\columnwidth}%
    \hbox\bgroup\kern-.2pt\vrule width.2pt%
      \vbox\bgroup\parindent0pt\relax
%    The 3pt is to cancel the -\lineskip from \displ@y
    \abovedisplayskip3pt \abovedisplayshortskip\abovedisplayskip
    \belowdisplayskip0pt \belowdisplayshortskip\belowdisplayskip
    \noindent}
}{%
  \par
%    Ensure that a lonely \[\] structure doesn't take up width less than
%    \hsize.
  \hrule height0pt width\hsize
  \egroup\vrule width.2pt\kern-.2pt\egroup
  \end{minipage}%
  \par\addvspace\medskipamount
}

%    Errors in output shouldn't be \texttt, and they may also be longer
%    than one line, but they should be identified by the "errorbullet".
%    (They were formerly labeled as aubsections.)
\newenvironment{erroro}[1]{%
  \par\addvspace\medskipamount
  \noindent\hangindent\parindent
  \hbox to\parindent{\errorbullet\hfil}\ignorespaces
  #1\par\vspace{\smallskipamount}
}{%
  \par\addvspace\medskipamount
}

%    The chapters are so short, perhaps we shouldn't call them by the
%    name `Chapter'. We make \chaptername read an argument in order to
%    remove a following \space or "{} " (both possibilities are present
%    in book.cls).

\renewcommand{\chaptername}[1]{}
\newcommand{\chapnum}[1]{\mdash #1\mdash }
\makeatletter
\def\@makechapterhead#1{%
  \vspace{1.5\baselineskip}%
  {\parindent \z@ \raggedright \reset@font
    \ifnum \c@secnumdepth >\m@ne
      \large\bfseries \chapnum{\thechapter}%
      \par\nobreak
      \vskip.5\baselineskip\relax
    \fi
    #1\par\nobreak
    \vskip\baselineskip
  }}
\makeatother

%    A command for ragged-right parbox in a tabular.
\newcommand{\rp}{\let\PBS\\\raggedright\let\\\PBS}

%    Non-indexed file name
\newcommand{\nfn}[1]{\texttt{#1}}

%    For the examples in the math spacing table.
%%\newcommand{\lspx}{\mbox{\rule{5pt}{.6pt}\rule{.6pt}{6pt}}}
%%\newcommand{\rspx}{\mbox{\rule[-1pt]{.6pt}{7pt}%
%%  \rule[-1pt]{5pt}{.6pt}}}
\newcommand{\lspx}{\mathord{\Rightarrow\mkern-1mu}}
\newcommand{\rspx}{\mathord{\mkern-1mu\Leftarrow}}
\newcommand{\spx}[1]{$\lspx #1\rspx$}

%    For a list of characters representing document input.
\newcommand{\clist}[1]{%
  \mbox{\ntt\spaceskip.2em plus.1em \xspaceskip\spaceskip#1}}

%    Fix weird \latex/ definition of rightmark.
\makeatletter
\def\rightmark{\expandafter\@rightmark\botmark{}{}}
%    Also turn off section marks.
\let\sectionmark\@gobble
\renewcommand{\chaptermark}[1]{%
  \uppercase{\markboth{\rhcn#1}{\rhcn#1}}}
\newcommand{\rhcn}{\thechapter. }
\makeatother

%    Include down to \section but not \subsection, in toc:
\setcounter{tocdepth}{1}

%    Subheadings for the bibliography
\newif\iffirstbibsubhead \firstbibsubheadtrue
\newcommand{\bibsubhead}[1]{%
  \iffirstbibsubhead \firstbibsubheadfalse
  \else \addvspace{\medskipamount}
  \fi
  \item[]\hspace*{-\leftmargin}\textbf{#1}\par
  \vspace{\smallskipamount}
}

\DeclareMathOperator{\ix}{ix}
\DeclareMathOperator{\nul}{nul}
\DeclareMathOperator{\End}{End}
\DeclareMathOperator{\xxx}{xxx}

\usepackage{hyperref}
\hypersetup{pdfborder=0 0 0}
\makeatletter
\let\oldcs\cs
\def\cs#1{\texorpdfstring{\oldcs{#1}}{\@backslashchar\@backslashchar#1}}
\let\cn\cs
\makeatother
\begin{document}
\chapter{Miscellaneous mathematical features}

\section{Matrices}\label{ss:matrix}

The \nipkg{amsmath} package provides some environments for
matrices\index{matrices} beyond the basic \env{array} environment of
\latex/. The \env{pmatrix}, \env{bmatrix}, \env{Bmatrix}, \env{vmatrix}
and \env{Vmatrix} have (respectively) $(\,)$, $[\,]$, $\lbrace\,\rbrace$,
$\lvert\,\rvert$, and $\lVert\,\rVert$ delimiters\index{delimiters}
built in. For naming consistency there is a
\env{matrix} environment sans delimiters. This is not entirely redundant
with the \env{array} environment; the matrix environments all use more
economical horizontal spacing than the rather prodigal spacing of the
\env{array} environment. Also, unlike the \env{array} environment, you
don't have to give column specifications for any of the matrix
environments; by default you can have up to 10 centered columns.%
\footnote{%%%%%%%%%%%%%%%%%%%%%%%%%%%%%%%%%%%%%%%%%%%%%%%%%%%%%%%%%%%%%%
More precisely: The maximum number of columns in a matrix is determined
by the counter |MaxMatrixCols| (normal value = 10), which you can change
if necessary using \latex/'s \cn{setcounter} or \cn{addtocounter}
commands.%
}\space%%%%%%%%%%%%%%%%%%%%%%%%%%%%%%%%%%%%%%%%%%%%%%%%%%%%%%%%%%%%%%%%%
(If you need left or right alignment in a column or other special
formats you may use \env{array}, or the \pkg{mathtools}
package which provides \verb+*+ variants%
\index{matrix@\indexenvironmentfont{matrix} environment!\verb+*+ variants}
of these environments with
an optional argument to specify left or right alignment.)

To produce a small matrix suitable for use in text, there is a
\env{smallmatrix} environment (e.g.,
\begin{math}
\bigl( \begin{smallmatrix}
  a&b\\ c&d
\end{smallmatrix} \bigr)
\end{math})
that comes closer to fitting within a single text line than a normal
matrix. Delimiters\index{delimiters} must be provided. (The
\pkg{mathtools} package provides |p|,|b|,|B|,|v|,|V| versions of
\env{smallmatrix},%
\index{smallmatrix@\indexenvironmentfont{smallmatrix} environment!variant delimiters}
as well as |*| variants%
\index{smallmatrix@\indexenvironmentfont{smallmatrix} environment!\verb+*+ variants}
as described above.)
The above example was produced by
\begin{verbatim}
\bigl( \begin{smallmatrix}
  a&b\\ c&d
\end{smallmatrix} \bigr)
\end{verbatim}

\cn{hdotsfor}|{|\<number>|}| produces a row of dots in a
matrix\index{matrices!ellipsis dots}\index{ellipsis dots!in
matrices}\index{dots|see{ellipsis dots}} spanning the given number of
columns. For example,
\begin{center}
\begin{minipage}{.3\columnwidth}
\noindent$\begin{matrix} a&b&c&d\\
e&\hdotsfor{3} \end{matrix}$
\end{minipage}%
\qquad
\begin{minipage}{.45\columnwidth}
\begin{verbatim}
\begin{matrix} a&b&c&d\\
e&\hdotsfor{3} \end{matrix}
\end{verbatim}
\end{minipage}%
\end{center}

The spacing of the dots can be varied through use of a square-bracket
option,\index{options!behavior of particular options}
for example, |\hdotsfor[1.5]{3}|.  The number in square brackets
will be used as a multiplier (i.e., the normal value is 1.0).
\begin{equation}\label{eq:D}
\begin{pmatrix} D_1t&-a_{12}t_2&\dots&-a_{1n}t_n\\
-a_{21}t_1&D_2t&\dots&-a_{2n}t_n\\
\hdotsfor[2]{4}\\
-a_{n1}t_1&-a_{n2}t_2&\dots&D_nt\end{pmatrix},
\end{equation}
\begin{verbatim}
\begin{pmatrix} D_1t&-a_{12}t_2&\dots&-a_{1n}t_n\\
-a_{21}t_1&D_2t&\dots&-a_{2n}t_n\\
\hdotsfor[2]{4}\\
-a_{n1}t_1&-a_{n2}t_2&\dots&D_nt\end{pmatrix}
\end{verbatim}

%% ------------------------------------------------------------------ %%

\section{Math spacing commands}

The \nipkg{amsmath} package slightly extends the set of math
spacing\index{horizontal spacing} commands, as shown below.
Both the spelled-out and abbreviated forms of these commands are robust,
and they can also be used outside of math.
\begin{ctab}{lll|lll}
Abbrev.& Spelled out& Example & Abbrev.& Spelled out& Example\\
\hline
\vstrut{2.5ex}
& no space& \spx{}& & no space & \spx{}\\
\cn{\,}& \cn{thinspace}& \spx{\,}&
  \altcnbang& \cn{negthinspace}& \spx{\!}\\
\cn{\:}& \cn{medspace}& \spx{\:}&
  & \cn{negmedspace}& \spx{\negmedspace}\\
\cn{\;}& \cn{thickspace}& \spx{\;}&
  & \cn{negthickspace}& \spx{\negthickspace}\\
& \cn{quad}& \spx{\quad}\\
& \cn{qquad}& \spx{\qquad}
\end{ctab}
For the greatest possible control over math spacing, use \cn{mspace}
and `math units'. One math unit, or |mu|, is equal to 1/18 em. Thus to
get a negative \cn{quad} you could write |\mspace{-18.0mu}|.

%% ------------------------------------------------------------------ %%

\section{Dots}

For preferred placement of ellipsis dots (raised or on-line) in various
contexts there is no general consensus. It may therefore be considered a
matter of taste. By using the semantically oriented commands
\begin{itemize}
\item \cn{dotsc} for \qq{dots with commas}
\item \cn{dotsb} for \qq{dots with binary operators/relations}
\item \cn{dotsm} for \qq{multiplication dots}
\item \cn{dotsi} for \qq{dots with integrals}
\item \cn{dotso} for \qq{other dots} (none of the above)
\end{itemize}
instead of \cn{ldots} and \cn{cdots}, you make it possible for your
document to be adapted to different conventions on the fly, in case (for
example) you have to submit it to a publisher who insists on following
house tradition in this respect. The default treatment for the various
kinds follows American Mathematical Society conventions:
\begin{center}
\begin{tabular}{@{}l@{}l@{}}
\begin{minipage}[t]{.54\textwidth}
\begin{verbatim}
Then we have the series $A_1, A_2,
\dotsc$, the regional sum $A_1
+A_2 +\dotsb $, the orthogonal
product $A_1 A_2 \dotsm $, and
the infinite integral
\[\int_{A_1}\int_{A_2}\dotsi\].
\end{verbatim}
\end{minipage}
&
\begin{minipage}[t]{.45\textwidth}
\noindent
Then we have the series $A_1,A_2,\dotsc$,
the regional sum $A_1+A_2+\dotsb$,
the orthogonal product $A_1A_2\dotsm$,
and the infinite integral
\[\int_{A_1}\int_{A_2}\dotsi.\]
\end{minipage}
\end{tabular}
\end{center}
For most situations, the undifferentiated \cn{dots} can be used, and
\nipkg{amsmath} will output the most suitable form based on the
immediate context; if an inappropriate form results, it can be corrected
after examining the output.

%% ------------------------------------------------------------------ %%

\section{Nonbreaking dashes}

A command \cn{nobreakdash} is provided to suppress the possibility of a
line break\index{line break|(} after the following hyphen or dash. For
example, if you write `pages 1\ndash 9' as |pages 1\nobreakdash--9|
then a line break will never occur between the dash and the 9. You can
also use \cn{nobreakdash} to prevent undesirable hyphenations in
combinations like |$p$-adic|. For frequent use, it's advisable to make
abbreviations, e.g.,
\begin{verbatim}
\newcommand{\p}{$p$\nobreakdash}% for "\p-adic"
\newcommand{\Ndash}{\nobreakdash--}% for "pages 1\Ndash 9"
%    For "\n dimensional" ("n-dimensional"):
\newcommand{\n}[1]{$n$\nobreakdash-\hspace{0pt}}
\end{verbatim}
The last example shows how to prohibit a line break\index{line break|)}
after the hyphen but allow normal hyphenation in the following
word. (It suffices to add a zero-width space after the hyphen.)

%% ------------------------------------------------------------------ %%

\section{Accents in math}

In ordinary \LaTeX{} the placement of the second accent in doubled math
accents is often poor. With the \nipkg{amsmath} package you
will get improved placement of the second accent:
$\hat{\hat{A}}$ (\cn{hat}|{\hat{A}}|).

The commands \cn{dddot} and \cn{ddddot} are available to produce triple
and quadruple dot accents in addition to the \cn{dot} and \cn{ddot}
accents already available in \latex/.

To get a superscripted hat or tilde character, load the \pkg{amsxtra}
package and use \cn{sphat} or \cn{sptilde}. Usage is \verb'A\sphat'
(note the absence of the \verb'^' character).


To place an arbitrary symbol in math accent position, or to get under
accents, see the \pkg{accents} package by Javier Bezos. (\nipkg{amsmath}
must be loaded before \pkg{accents}.)

%% ------------------------------------------------------------------ %%

\section{Roots}

In ordinary \latex/ the placement of root indices is sometimes not so
good: $\sqrt[\beta]{k}$ (|\sqrt||[\beta]{k}|).  In the
\nipkg{amsmath} package \cn{leftroot} and \cn{uproot} allow you to adjust
the position of the root:%
\index{options!adjust positioning}
\begin{verbatim}
  \sqrt[\leftroot{-2}\uproot{2}\beta]{k}
\end{verbatim}
will move the beta up and to the right:
$\sqrt[\leftroot{-2}\uproot{2}\beta]{k}$. The negative argument used
with \cn{leftroot} moves the $\beta$ to the right. The units are a small
amount that is a useful size for such adjustments.

%% ------------------------------------------------------------------ %%

\section{Boxed formulas}

The command \cn{boxed} puts a box around its
argument, like \cn{fbox} except that the contents are in math mode:
\begin{equation}
\boxed{\eta \leq C(\delta(\eta) +\Lambda_M(0,\delta))}
\end{equation}
\begin{verbatim}
  \boxed{\eta \leq C(\delta(\eta) +\Lambda_M(0,\delta))}
\end{verbatim}

%% ------------------------------------------------------------------ %%

\section{Over and under arrows}

Basic \latex/ provides \cn{overrightarrow} and \cn{overleftarrow}
commands. Some additional over and under arrow commands are provided
by the \nipkg{amsmath} package to extend the set:

\begin{tabbing}
\qquad\=\ncn{overleftrightarrow}\qquad\=\kill
\> \cn{overleftarrow}    \> \cn{underleftarrow} \+\\
   \cn{overrightarrow}    \> \cn{underrightarrow} \\
   \cn{overleftrightarrow}\> \cn{underleftrightarrow}
\end{tabbing}

%% ------------------------------------------------------------------ %%

\section{Extensible arrows}

\cn{xleftarrow} and \cn{xrightarrow} produce
arrows\index{arrows!extensible} that extend automatically to accommodate
unusually wide subscripts or superscripts. These commands take one
optional argument\index{options!behavior of particular options}
(the subscript) and one mandatory argument (the
superscript, possibly empty):
\begin{equation}
A\xleftarrow{n+\mu-1}B \xrightarrow[T]{n\pm i-1}C
\end{equation}
\begin{verbatim}
  \xleftarrow{n+\mu-1}\quad \xrightarrow[T]{n\pm i-1}
\end{verbatim}

%% ------------------------------------------------------------------ %%

\section{Affixing symbols to other symbols}

\latex/ provides \cn{stackrel} for placing a
superscript\index{subscripts and superscripts} above a binary relation.
In the \nipkg{amsmath} package there are somewhat more general commands,
\cn{overset} and \cn{underset}, that can be used to place one symbol
above or below another symbol, whether it's a relation or something
else. The input |\overset{*}{X}| will place a superscript-size {\small$*$} above
the {\small$X$}: {\small$\smash{\overset{*}{X}}$};
\cn{underset} is the analog for adding a
symbol underneath. The command \cn{overunderset} is a combination of
these, taking three arguments to place superscript sized expressions
above and below the same base.

See also the description of \cn{sideset} in \secref{sideset}.

%% ------------------------------------------------------------------ %%

\section{Fractions and related constructions}

\subsection{The \cn{frac}, \cn{dfrac}, and \cn{tfrac} commands}

The \cn{frac} command, which is in the basic command set of
\latex/,\index{fractions|(} takes two arguments\mdash numerator and
denominator\mdash and typesets them in normal fraction form. The
\nipkg{amsmath} package provides also \cn{dfrac} and \cn{tfrac} as
convenient abbreviations for |{\displaystyle\frac| |...| |}|
and\indexcs{textstyle}\relax
\indexcs{displaystyle} |{\textstyle\frac| |...| |}|.
\begin{equation}
\frac{1}{k}\log_2 c(f)\quad\tfrac{1}{k}\log_2 c(f)\quad
\sqrt{\frac{1}{k}\log_2 c(f)}\quad\sqrt{\dfrac{1}{k}\log_2 c(f)}
\end{equation}
\begin{verbatim}
\begin{equation}
\frac{1}{k}\log_2 c(f)\;\tfrac{1}{k}\log_2 c(f)\;
\sqrt{\frac{1}{k}\log_2 c(f)}\;\sqrt{\dfrac{1}{k}\log_2 c(f)}
\end{equation}
\end{verbatim}

\subsection{The \cn{binom}, \cn{dbinom}, and \cn{tbinom} commands}

For binomial expressions\index{binomials} such as $\binom{n}{k}$
\nipkg{amsmath} has \cn{binom}, \cn{dbinom} and \cn{tbinom}:
\begin{equation}
2^k-\binom{k}{1}2^{k-1}+\binom{k}{2}2^{k-2}
\end{equation}
\begin{verbatim}
2^k-\binom{k}{1}2^{k-1}+\binom{k}{2}2^{k-2}
\end{verbatim}

\subsection{The \cn{genfrac} command}

The capabilities of \cn{frac}, \cn{binom}, and their variants are
subsumed by a generalized fraction command \cn{genfrac} with six
arguments. The last two correspond to \cn{frac}'s numerator and
denominator; the first two are optional delimiters\index{delimiters} (as
seen in \cn{binom}); the third is a line thickness override (\cn{binom}
uses this to set the fraction line thickness to 0\mdash i.e., invisible);
and the fourth argument is a mathstyle override: integer values 0\ndash 3
select respectively \cn{displaystyle}, \cn{textstyle}, \cn{scriptstyle},
and \cn{scriptscriptstyle}. If the third argument is left empty, the line
thickness defaults to `normal'.

\begin{center}\begin{minipage}{.85\columnwidth}
\raggedright \normalfont\ttfamily \exhyphenpenalty10000
\newcommand{\ma}[1]{%
  \string{{\normalfont\itshape#1}\string}\penalty9999 \ignorespaces}
\string\genfrac \ma{left-delim} \ma{right-delim} \ma{thickness}
\ma{mathstyle} \ma{numerator} \ma{denominator}
\end{minipage}\end{center}
To illustrate, here is how \cn{frac}, \cn{tfrac}, and
\cn{binom} might be defined.
\begin{verbatim}
\newcommand{\frac}[2]{\genfrac{}{}{}{}{#1}{#2}}
\newcommand{\tfrac}[2]{\genfrac{}{}{}{1}{#1}{#2}}
\newcommand{\binom}[2]{\genfrac{(}{)}{0pt}{}{#1}{#2}}
\end{verbatim}
If you find yourself repeatedly using \cn{genfrac} throughout a document
for a particular notation, you will do yourself a favor (and your
publisher) if you define a meaningfully-named abbreviation for that
notation, along the lines of \cn{frac} and \cn{binom}.

The primitive generalized fraction commands \cs{over}, \cs{overwithdelims},
\cs{atop}, \cs{atopwithdelims}, \cs{above}, \cs{abovewithdelims} produce
warning  messages if used with the \nipkg{amsmath} package, for reasons
discussed in \fn{technote.tex}.

%% ------------------------------------------------------------------ %%

\end{document}
